\section{Discussion}
\label{sec:discussion}

\subsection{What the skill result does and does not show}

The system exceeds a four-predictor observational benchmark by $0.054$ ACC over leads 1--18
(Sec.~\ref{sec:skill}). Three qualifications are essential to reading that number correctly.

First, the benchmark is deliberately strong. It is an ordinary least-squares map from Ni\~no
3.4, warm-water volume, and their three-month tendencies -- that is, from the state variables
that recharge-oscillator theory identifies as the relevant ones
\citep{Jin1997} -- onto the same latent space, decoded through the same
frozen decoder. The comparison therefore isolates what the learned operator contributes beyond
a linear reading of the recognised predictors, rather than measuring skill against a
deliberately weak reference.

Second, the advantage is narrow rather than uniform. It is significant for El Ni\~no events
but not La Ni\~na (Sec.~\ref{sec:event-auc}), for boreal-winter targets but not summer
(Sec.~\ref{sec:seasonality}), and per-lead only at leads 1--7 and 12--16. Sections
\ref{sec:event-auc} and \ref{sec:seasonality} describe one effect, not two: the operator's
residual contributes to the mature winter phase of warm events and contributes essentially
nothing to the summer transition or to cold events. The sign symmetry of a linear operator
offers a natural explanation for the warm/cold asymmetry, and the seasonal structure follows
ENSO's own phase-locking, but we note that these are consistent interpretations rather than
tests: neither is established by the results presented here.

Third, the most quotable statistic in the paper is not a result. The ACC$=0.5$ crossing moves
from 9.8 to 12.4 months, but its 95\% interval is $[-0.33, +9.93]$ months and does not exclude
zero. Threshold-crossing statistics on noisy curves are unstable by construction, and we report
this one descriptively only.

\subsection{How much interpretability the operator actually provides}

All 40 trained models converge on the same three periods, and the ordering of their damping
rates separates from an identical initialisation -- the strongest emergent result in
Sec.~\ref{sec:modal-spectrum}, precisely because the damping rates, unlike the periods, begin
undifferentiated. The 2.41-year quasi-biennial mode becoming the least damped is a physically
meaningful outcome that a black-box forecaster would not expose.

The interpretability claim should nonetheless be bounded carefully. The periods are initialised
in-band and are therefore displaced from, rather than discovered independently of, their
starting values. Stability is imposed by the parameterisation, so the fact that all modes are
damped carries no information. Only the leading mode is dynamically excited above an isotropic
baseline; the 5.73-year mode fails three independent tests and should not be presented as an
established ENSO band. And the ablation in Sec.~\ref{sec:ablations} shows that skill is
statistically indistinguishable across one to four oscillators. Taken together, the honest
statement is that the operator recovers \emph{one} dominant, physically recognisable mode plus
damped memory, and that the specific three-oscillator structure is a modelling convenience
rather than a result.

\subsection{Three null results, and what they close}

Sections~\ref{sec:ablations} reports three controlled ablations, none of which improved skill:
restricting the operator to an ENSO-correlated subspace \emph{hurt} it ($-0.068$ ACC), adding a
fourth oscillator did not survive cross-fold confirmation, and an explicit amplitude
nonlinearity produced no measurable change.

The first is informative rather than merely negative. It resolves an apparent contradiction
with Sec.~\ref{sec:latent-coords}: ENSO's forecast \emph{output} is concentrated in a single
latent coordinate, but the \emph{state} required to predict that coordinate is not, which is
what recharge-oscillator physics predicts if subsurface structure rather than surface
temperature carries the memory.

The second is methodological. A single-fold development experiment scored the fourth oscillator
at $t=5.93$ with 6 of 6 seeds improving; the four-fold paired test scored the same
configuration at $t=1.31$ with 20 of 40. The effect survived only on the fold that had selected
it. We report this because the failure mode is general and easy to fall into: with per-seed
noise of $0.005$ and a development fold of 99 verification windows, a single-fold experiment
can produce a decisively significant result for a configuration that does not generalise.

The third motivated the analysis in Sec.~\ref{sec:amplitude-tradeoff}. Given the capacity to
amplify -- a bounded self-amplification term permitting up to roughly a doubling over 24
months -- the operator instead learned additional damping in five of six oscillator instances.

\subsection{Amplitude is a calibration choice, not a dynamical limit}

We consider this the most transferable result in the paper. The forecast underestimates event
amplitude substantially: a regression slope of $0.516$ at lead 6
(Sec.~\ref{sec:case-studies}). The natural reading is that a damped linear operator cannot grow
an anomaly. That reading is wrong on both counts. The operator can amplify and declines to
(Sec.~\ref{sec:ablations}), and the unblended operator output already carries a slope of
$0.729$ -- more than the blended forecast retains. The shortfall is produced by the blend
weight, which is set to maximise correlation skill, and full amplitude is recoverable at a
measured, exactly computable cost of $0.19$ ACC (Sec.~\ref{sec:amplitude-tradeoff}).

This decomposition is not specific to our architecture. Any forecast system that blends towards
a conservative reference, or that is trained or tuned against a correlation-like score, faces
the same exchange: correlation is invariant to a uniform rescaling of the forecast, so
optimising it provides no pressure towards calibrated amplitude, and shrinkage toward the mean
reduces squared error whenever the forecast is uncertain. Systems tuned this way should be
expected to be systematically under-dispersive, and the magnitude of that effect is measurable
without retraining wherever the unblended forecast has been retained. We suggest reporting an
amplitude slope alongside correlation as a matter of routine.

\subsection{Limitations}

The most consequential limitation is the absence of an external baseline. All comparisons here
are internal -- against the observational benchmark, persistence, and climatology -- so this
paper establishes what the architecture contributes relative to a linear reading of its own
predictors, and does \emph{not} establish where it stands relative to operational dynamical
forecasts or to previously published learned ENSO models. That comparison is the natural next
step and is required before any claim of competitive skill.

Skill is strongly regime-dependent. Per-fold mean ACC over leads 1--18 ranges from $0.784$
(1981--1991) to $0.336$ (2002--2013), and the quiet 2002--2013 decade is also roughly twice as
seed-sensitive as the others. A single pooled skill number understates this spread
considerably.

The effective sample is small in the way that matters. The record contains 43 years but only 11
El Ni\~no and 14 La Ni\~na episodes exceeding $1.0\deg$C, which is why event-discrimination
intervals are wide and why the paired year-block bootstrap, rather than window-level
resampling, is used throughout. Related, the event threshold used here is $\pm1.0\deg$C rather
than the conventional $\pm0.5$; at the conventional threshold the El Ni\~no advantage is not
statistically resolvable.

Finally, leads 19--24 are identical to the benchmark by construction, and forecasts of
boreal-summer targets at long lead show no demonstrable skill.

\section{Conclusions}
\label{sec:conclusions}

We have presented an ENSO forecast system in which representation learning and dynamics are
separated: a transformer autoencoder with circulation-shaped attention windows learns a
32-dimensional state, and a seasonally modulated linear operator built from damped oscillators
advances it. The system reaches ACC $0.619$ over leads 1--18, exceeding a strong
observational-predictor benchmark by $0.054$ ($[0.003, 0.104]$), with the advantage
concentrated in El Ni\~no events and boreal-winter targets at leads 12--17.

Because the dynamics are structurally simple, they can be inspected. All 40 trained models
converge on a dominant quasi-biennial mode at $2.41 \pm 0.04$ years which becomes the least
damped of the three from an initialisation where the damping rates are identical. We have been
explicit about the limits of that interpretability: the periods are initialised in-band,
stability is imposed, only the leading mode is excited, and skill is insensitive to whether
one, two, three, or four oscillators are used.

Three controlled ablations returned null or negative results, and we report them because they
constrain the design: the full 32-dimensional state is necessary despite ENSO's output being
concentrated in one coordinate; a fourth oscillator does not survive cross-fold confirmation
despite appearing decisive on a single development fold; and an explicit amplitude
nonlinearity changes nothing.

The amplitude finding is the one we expect to generalise beyond this system. The forecast's
underestimation of event amplitude is not a limitation of its dynamics but a consequence of
optimising a correlation-like objective, and the exchange rate between amplitude and
correlation can be computed exactly, without retraining, from forecasts that have already been
made. We recommend that forecast systems of this kind report an amplitude slope alongside
correlation, and that the unblended forecast be retained so the trade remains measurable.

The immediate priority for further work is comparison against established dynamical and learned
ENSO forecasts on a common verification period. Beyond that, given three null results from
architectural modification, the more promising direction is additional training data --
a longer reanalysis record or pre-training on coupled-model output -- rather than further
elaboration of the operator.
